% ===================================================================
%  TÁBLA Uimh. 16 — An Siopa Mór. — An Basár. — Na Bréagáin.
%  Traduction 2026 d'apres texto/io, controlee sur texto/fr, texto/en
%  et texto/es. Consigne : voir texto/ga/10-tabla-01.tex.
%
%  Le tableau d'astronomie (bloc 15-1) porte des renvois a LETTRES,
%  (a) a (f), graves sur la planche elle-meme : ils se rangent sous le
%  numero 85, comme le dit gravuri/literi.json. Ce bloc est aussi le
%  troisieme ecart declare de renvoji.py.
%
%  LA MENAGERIE DE CARTON EST UN CAS A PART, et c'est le dernier
%  enseignement lexical de la colonne. Les betes de la ferme et du
%  bois irlandais ont toutes leur mot — « sionnach » le renard,
%  « mac tíre » le loup, « luch » la souris, « francach » le rat,
%  « ialtóg » la chauve-souris. Les betes qu'aucun Irlandais n'a
%  jamais vues vivantes arrivent avec le livre d'images :
%  « sioráf », « hipapotamas », « dromadaire », « leopard »,
%  « hiéna », « crogall », « pantar », « strus ». La frontiere ne
%  passe pas entre l'ancien et le nouveau ni entre le dedans et le
%  dehors, mais entre CE QU'ON A VU ET CE QU'ON A LU. C'est encore le
%  lieu de l'apprentissage, pousse a sa limite : la ou l'objet ne
%  vient jamais, seul son nom voyage.
%
%  DERNIER TABLEAU DE LA COLONNE IRLANDAISE.
% ===================================================================

%%K t16-apar-1 apar
\VUsaut{4.0mm}
\VUtitre{15.0pt}{60}{TÁBLA Uimh. 16}
\VUsaut{2.0mm}
\VUpk{13.2pt}{40}{An Siopa Mór. --- An Basár.}
\VUpk{13.2pt}{40}{Na Bréagáin.}
\VUsaut{3.4mm}

%%K t16-01-1 p
1. --- Tá an bhliain nua ag druidim linn. Tá na siopaí móra tar éis a dtaispeántais
\VUgras{bréagán} \textsuperscript{(1)} agus bronntanas na bliana nua a ullmhú.

%%K t16-01-2 p
D'iarr mé ar mo sheanathair mé a thabhairt chun an bhasáir chun an taispeántas breá
sin a fheiceáil, agus thoiligh sé leis.

%%K t16-02-1 p
2. --- Ar dtús stadamar os comhair sciatha breátha arm, ina bhfuil \VUgras{raidhfil,}
\VUgras{céipí, sabar, crios claímh} agus péire \VUgras{gualainníní,} nó
\VUgras{clogad,} \VUgras{cathéide} \textsuperscript{(3)} agus \VUgras{piostal}
\textsuperscript{(4)}. Chuir sin ar fad i bhfad níos mó suime ionam ná na
\VUgras{teilgeoirí solais} \textsuperscript{(2)}.

%%K t16-tit-1 sub
\VUsaut{3.0mm}
\VUpk{10.2pt}{40}{Radharcanna breátha.}
\VUsaut{2.6mm}

%%K t16-03-1 p
3. --- San urrann chéanna bhí \VUgras{fear faire} \textsuperscript{(5)} ina
\VUgras{bhosca faire}; ba dhóigh leat go raibh sé ag faire os comhair gairdín
ainmhithe cairtchláir ar fad. Nach iomaí ainmhí a bhí ann: \VUgras{Pearóid}
\textsuperscript{(6)} ar a fara, \VUgras{srónbheannach} \textsuperscript{(7)} agus
beann ar a shrón, \VUgras{réinfhia} \textsuperscript{(8)}, \VUgras{strus}
\textsuperscript{(9)}, \VUgras{moncaí} \textsuperscript{(10)} ag briseadh cnó,
\VUgras{sionnach} \textsuperscript{(11)} ag breith cearc leis, \VUgras{crogall}
\textsuperscript{(12)} agus a chraos ollmhór ar leathadh, \VUgras{camall}
\textsuperscript{(13)}, \VUgras{cú caol} galánta \textsuperscript{(14)},
\VUgras{pantar} \textsuperscript{(15)}, ansin dhá ainmhí nach raibh aithne agam orthu
agus arb iad, de réir mar a dúradh, \VUgras{easóg} \textsuperscript{(16)} agus
\VUgras{hiéna} \textsuperscript{(17)} iad. Bhí \VUgras{liopard}
\textsuperscript{(18)} agus \VUgras{mac tíre} \textsuperscript{(21)} ann freisin;
díreach in aice leo bhí \VUgras{francaigh} bheaga ghleoite \textsuperscript{(19)}
agus \VUgras{lucha} bídeacha \textsuperscript{(20)} rubair. Ar deireadh chuir
\VUgras{hipapotamas} \textsuperscript{(22)} agus \VUgras{dromadaire}
\textsuperscript{(23)} dronnach clabhsúr ar an mbailiúchán. D'fhanfainn ann uaireanta
níos faide murach go raibh in aice láimhe machaire iontach lán de \VUgras{shaighdiúirí
luaidhe} \textsuperscript{(29)} a tharraing m'aird.

%%K t16-tit-2 sub
\VUsaut{4.0mm}
\VUpk{10.2pt}{40}{Saighdiúirí agus cogadh.}
\VUsaut{2.8mm}

%%K t16-04-1 p
4. --- Is mór le mo sheanathair, ar \VUgras{easlán} \textsuperscript{(24)} é agus a
raibh air an t-arm a fhágáil de bharr \VUgras{créachta} a thuill an
\VUgras{bonn míleata} dó, cur síos a dhéanamh dom ar na \VUgras{feachtais} inar
ghlac sé páirt. Bhí áthas air agus é ag míniú ghluaiseachtaí éagsúla an airm bhig
mhionsamhlaigh. Stad grúpa fíorshaighdiúirí a bhí ar cuairt ar an mbasár in aice
linn chun éisteacht leis na mínithe: \VUgras{saighdiúir innealtóireachta}
\textsuperscript{(25)}, \VUgras{sealgaire Alpach} \textsuperscript{(26)} agus
\VUgras{bairéad} ar a cheann, \VUgras{sáirsint mór} \textsuperscript{(27)} coisithe
agus \VUgras{sáirsint airtléire} \textsuperscript{(28)} agus \VUgras{galún} óir ar a
mhuinchille.

%%K t16-05-1 p
5. --- Bhí mo sheanathair chomh bródúil sin go raibh lucht éisteachta chomh
taibhseach sin aige, gur chum sé as a cheann plean iomlán catha.

%%K t16-05-2 p
Bhí an chathair dhaingnithe, agus a \VUgras{rampair} agus a \VUgras{díoga} aici, faoi
\VUgras{léigear} ag an \VUgras{arm namhad,} arm a bhí réidh chun \VUgras{ionsaí} a
dhéanamh tar éis \VUgras{bombardaithe} fhada. Ach bhain \VUgras{na daoine faoi
léigear}, a raibh an t-ocras ag brú orthu, triail as \VUgras{ruathar} a thabhairt ar
\VUgras{champa} lucht an léigir. Tar éis dóibh an \VUgras{t-ionsaí} a chur ar gcúl
le \VUgras{comhrac} dianseasmhach timpeall an \VUgras{champa}, chuir lucht an
léigir \VUgras{buacach} a gcuid \VUgras{marcshluaite} amach chun tóir a chur ar na
hionsaitheoirí \VUgras{cloíte}. \VUgras{Chúlaigh} siad sin agus
\VUgras{mairbh} agus \VUgras{lucht créachta} fágtha acu ar fud an \VUgras{mhachaire
chatha}. Thug \VUgras{lucht na sínteán} an dream deireanach sin leo isteach san
\VUgras{otharcharr} le go dtabharfadh an \VUgras{máinlia} aire dóibh.

%%K t16-05-3 p
In ainneoin frithsheasamh laochta roinnt \VUgras{cathlán} a rinne \VUgras{cearnóg} díobh
féin, ba iomlán é an \VUgras{bua,} agus \VUgras{theith} an chuid eile den arm agus
a lán \VUgras{príosúnach} fágtha ina ndiaidh acu. Chuaigh an namhaid chun a
\VUgras{bhua} a chur i gcrích trí sheilbh a ghlacadh ar an gcathair. B'fhéidir gurb é
sin deireadh an fheachtais, agus tar éis \VUgras{sos cogaidh} go bhféadfaí
\VUgras{conradh síochána} a shíniú.

%%K t16-tit-3 sub
\VUsaut{4.4mm}
\VUpk{10.2pt}{40}{Bronntanais na bliana nua.}
\VUsaut{2.8mm}

%%K t16-06-1 p
6. --- D'iarr na saighdiúirí óga, a bhí ag gáire agus iad ag éisteacht le hinsint
phictiúrtha an tseanshaighdiúra, roinnt mínithe eile air. Maidir liom féin, chuaigh
mé timpeall an tsiopa chun féachaint ar \VUgras{chrosbhogha} breá
\textsuperscript{(32)} agus ar \VUgras{bhogha} \textsuperscript{(33)} agus a
\VUgras{shaighead} \textsuperscript{(31)} leis. Fuair mé ansin bailiúchán eile
ainmhithe: dhá \VUgras{éan cheoil} \textsuperscript{(34)}, \VUgras{filiméala}
uathoibríoch agus \VUgras{ceolaire} ag canadh in ard a ghutha, agus sraith ainmhithe
aisteacha: \VUgras{ialtóg} \textsuperscript{(35)}, \VUgras{bóa}
\textsuperscript{(36)}, \VUgras{turtar} \textsuperscript{(37)}, \VUgras{rón}
\textsuperscript{(38)}, \VUgras{míol mór} \textsuperscript{(39)} agus \VUgras{siorc}
\textsuperscript{(40)} de cheirt.

%%K t16-07-1 p
7. --- Níor stad mé i bhfad chun breathnú orthu, mar chonaic mé an rud a bhí i mo
bhrionglóidí: \VUgras{amharclann phuipéad} sárbhreá \textsuperscript{(41)} agus
\VUgras{maisiúcháin} iontacha agus \VUgras{cuirtín} dearg álainn uirthi. Ar an
\VUgras{stáitse} ba dhóigh leat go raibh beirt aisteoirí ag imirt \VUgras{páirte} i
\VUgras{ndráma} an-suimiúil, mar bhí an \VUgras{lucht féachana} beag cairtchláir a
bhí curtha sna loighsí nó ina suí ar na \VUgras{cathaoireacha móra} agus sa
\VUgras{phairtéar} taobh thiar de \VUgras{stiúrthóir na ceolfhoirne} ag bualadh bos
go beoga.

%%K t16-07-2 p
Faraor, bhí an amharclann sin an-daor.

%%K t16-08-1 p
8. --- Rud eile, ba dheacair na bréagáin a roghnú, mar bhí gach cineál bréagán
carntha sna fuinneoga: \VUgras{Airlicín} \textsuperscript{(42)},
\VUgras{Puinseoinéal} \textsuperscript{(43)}, \VUgras{Pierrot} \textsuperscript{(44)},
\VUgras{ceanna cait} \textsuperscript{(45)} ag croitheadh a sciathán, \VUgras{loin
dhubha} \textsuperscript{(46)} feadaíola, \VUgras{púdail} ag tafann,
\VUgras{fónagraf} \textsuperscript{(47)} agus \VUgras{bréagáin foighne.} Chuir ceann
de na bréagáin sin an-siamsa orm: bhí \VUgras{cath farraige}
\textsuperscript{(48)} á léiriú aige, cath ina raibh \VUgras{long} á hionsaí ag
\VUgras{foghlaithe mara} in aice le tír a raibh \VUgras{crainn chnó cócó} curtha
inti.

%%K t16-noto-1 noto
\VUnotes{143.2mm}{%
\textit{(*) Féach an tábla Uimh. 6.}%
}

%%K t16-09-1 p
9. --- D'éirigh liom breathnú go réidh ar na hiontais sin go léir, mar ní raibh ach
beagán daoine sa siopa, ar éigean cúpla fánaí, \VUgras{dragún} \textsuperscript{(49)}
agus \VUgras{bailitheoir} \textsuperscript{(50)} a chuir \VUgras{bille malairte} i
láthair ag an bpríomh-\VUgras{chiste} \textsuperscript{(51)}.

%%K t16-10-1 p
10. --- D'fhiafraigh mé de mo sheanathair cad chuige a n-úsáidtear na leabhair a bhí
curtha ar chláir taobh thiar de \VUgras{bhean an chiste}; d'fhreagair sé gur
\VUgras{leabhair chuntais} iad.

%%K t16-tit-4 sub
\VUsaut{3.6mm}
\VUpk{10.2pt}{40}{Ag stánadh timpeall an tsiopa.}
\VUsaut{2.8mm}

%%K t16-11-1 p
11. --- Ag fágáil urrann na mbréagán dúinn chuamar chuig an diallaitlann chun
súil a chaitheamh ar na \VUgras{diallaití} \textsuperscript{(53)}, na
\VUgras{stíoróipí} \textsuperscript{(54)}, na \VUgras{srianta}
\textsuperscript{(55)}, na \VUgras{béalbhaigh} \textsuperscript{(56)} agus na
\VUgras{fuipeanna.} Fad is a thug \VUgras{ceann na hurrainne}
\textsuperscript{(57)} roinnt eolais do mo sheanathair, chuaigh mé ag féachaint ar
thaispeántas na \VUgras{bpoirceallán} agus na \VUgras{gcriostal}
\textsuperscript{(58)} ina bhfuil \VUgras{babhlaí sailéid} breátha agus
\VUgras{taephotaí} mine \textsuperscript{(59)}.

%%K t16-12-1 p
12. --- Chun dul ar an gcéad urlár ghabhamar taobh thiar d'fhuinneog na
\VUgras{gcorsaí} \textsuperscript{(60)}, in aice le siopa na stocaí, mar a raibh
\VUgras{brístí} an-bhreátha \textsuperscript{(63)} ar taispeáint. In aice leis sin
chuir \VUgras{cailín siopa} \textsuperscript{(61)} ar \VUgras{cheannaitheoir mná}
\textsuperscript{(62)} \VUgras{fabraicí} síoda breátha \textsuperscript{(64)} a
roghnú.

%%K t16-12-2 p
Agus mé ag dul suas an staighre, thug mé faoi deara go bhfuil an siopa maisithe le
\VUgras{laindéir} Sheapánacha \textsuperscript{(65)}, le \VUgras{parasóil}
\textsuperscript{(66)} agus le \VUgras{crainn phailme} ollmhóra \textsuperscript{(70)}.

%%K t16-13-1 p
13 . --- Ar an gcéad urlár ní raibh mórán le feiceáil; ní raibh ann ach troscán (*),
\VUgras{taisceadáin iarainn} \textsuperscript{(67)}, \VUgras{cófraí tarraiceán}
\textsuperscript{(68)} agus \VUgras{pramanna} \textsuperscript{(69)}. Ach bhí
radharc iomlán an tsiopa an-\VUgras{taitneamhach.}

%%K t16-14-1 p
14. --- Ag ár gcosa chonaic mé na huirlisí ceoil: \VUgras{cláirseacha}
\textsuperscript{(71)}, \VUgras{giotáir} \textsuperscript{(72)},
\VUgras{maindilíní} \textsuperscript{(73)}, \VUgras{coirnéid comhla}
\textsuperscript{(74)}, \VUgras{cláirnéidíní} \textsuperscript{(75)}, veidhlíní agus
a \VUgras{mbogha} \textsuperscript{(76)} leo, agus fiú \VUgras{druma mór}
\textsuperscript{(77)}. Beagán níos faide uainn, ag urrann an pháipéir, bhí
\VUgras{mac léinn} \textsuperscript{(78)} ag margáil le \VUgras{cúntóir siopa}
\textsuperscript{(79)} faoi chás cóipleabhar. In aice leo d'aithin mé
\VUgras{aibítirleabhar} breá \textsuperscript{(80)}.

%%K t16-14-2 p
I lár an tsiopa bhí \VUgras{crann Nollag} ard \textsuperscript{(81)} curtha in airde
agus é lán de choinnle, de mhionrudaí agus de gach cineál bréagáin: \VUgras{capaill
adhmaid} \textsuperscript{(82)}, \VUgras{bábóga} \textsuperscript{(83)}, srl.

%%K t16-14-3 p
Scaip \VUgras{siopa na gcumhrán} \textsuperscript{(84)}, agus a chuid
\VUgras{taosanna fiacla} agus \VUgras{cumhrán} éagsúil ann, boladh an-bhreá a tháinig
inár dtreo.

%%K t16-tit-5 sub
\VUsaut{3.4mm}
\VUpk{10.2pt}{40}{Ceannaímid rud éigin.}
\VUsaut{2.8mm}

%%K t16-15-1 p
15. --- Ó thosaigh mo sheanathair ag ceapadh go raibh an t-am rófhada, chuamar síos
an staighre mór. Stad sé tamall os comhair \VUgras{tábla na réalteolaíochta}
\textsuperscript{(85)}, tábla a raibh air, mar a dúirt sé liom, \VUgras{cóiméid}
\textsuperscript{(a)}, \VUgras{urú gealaí agus urú gréine} \textsuperscript{(b)},
rós na ngaoth agus \VUgras{na ceithre hairde} \textsuperscript{(c)} \VUgras{(An
Tuaisceart, An Deisceart, An tOirthear agus An tIarthar)}, léarscáil an dá
\VUgras{leathsféar} \textsuperscript{(d)}, na \VUgras{pláinéid}
\textsuperscript{(e)} agus \VUgras{grianchóras} \textsuperscript{(f)}. Níor thuig mé
faic de sin ar fad, agus chuir sé i bhfad níos mó suime ionam an éide agus galúin uirthi a
bhí ar \VUgras{bhuachaill seachadta} \textsuperscript{(86)} a ghabh tharainn, agus
na \VUgras{feananna} breátha \textsuperscript{(87)} a bhí in aice liom, cois na
\VUgras{sparán airgid} \textsuperscript{(88)} agus na \VUgras{vaidhcín}
\textsuperscript{(89)}.

%%K t16-16-1 p
16. --- Bhí meas agam freisin ar an taispeántas ar \VUgras{chleití}
\textsuperscript{(90)} agus ar \VUgras{bhúclaí creasa} \textsuperscript{(91)}, agus ar
na \VUgras{bláthanna saorga} breátha \textsuperscript{(94)} a bhí curtha go
maisiúil i gciseáin. Bhí na \VUgras{sailchuacha}, na \VUgras{lilí sráide}, na
\VUgras{bú} \textsuperscript{(95)}, na \VUgras{geiréiniamaí} \textsuperscript{(96)},
na \VUgras{lilí} \textsuperscript{(97)} agus na \VUgras{caipisíní}
\textsuperscript{(98)} cóipeáilte go han-mhaith.

%%K t16-tit-6 sub
\VUsaut{4.4mm}
\VUpk{10.2pt}{40}{Bronntanas ó mo sheanathair.}
\VUsaut{2.8mm}

%%K t16-17-1 p
17. --- D'iarr mé ar mo sheanathair ceann de na \VUgras{scuaba fiacla}
\textsuperscript{(92)} a bhí i mbosca in aice le \VUgras{bioráin ghruaige}
\textsuperscript{(93)} a cheannach dom. Thoiligh sé leis agus chuaigh mé féin ag íoc
as mo cheannach ag an gciste, mar a raibh \VUgras{coinsíneacht} mhór
\textsuperscript{(100)} á hullmhú do \VUgras{chliant} \textsuperscript{(99)}.

%%K t16-18-1 p
18. --- Go tobann tharla corraíl bheag i measc na ndaoine a bhí stoptha in aice leis
na \VUgras{cré-umhaí} ealaíne \textsuperscript{(101)}. Bhí \VUgras{gadaí}
\textsuperscript{(102)} díreach gafa i mbun a choire ag \VUgras{cigire}
\textsuperscript{(103)}, agus é ag iarraidh duine a robáil. Cuireadh fios ar phóilín
chun é a \VUgras{ghabháil,} agus díreach agus sinn ag dul amach, bhí an coirpeach á
thabhairt chun an phríosúin.
